\section{Equilibrium and Log-Linearization}
\label{sec:equilibrium}

A sequential competitive equilibrium consists of household policy functions,
firm price and wage policies, and price sequences such that households
maximize \eqref{eq:hh_objective}--\eqref{eq:bc}, firms optimize
\eqref{eq:production} subject to monopsonistic labor supply
\eqref{eq:labor_supply} and Calvo pricing \eqref{eq:optimal_price}, and
all markets clear. Under monopsony, the labor market clears at
$w_t(\ell) = \mu_t^w(\ell) \cdot \mathrm{MPL}_t(\ell)$---below the
competitive wage---while goods and asset markets clear in the standard
fashion. The government runs a balanced budget with constant bond supply
$B_t = \bar B$.

\paragraph{Key log-linearized equations.}
The monopsony-adjusted Phillips curve (derived in Section~\ref{sec:results}):
\begin{equation}
  \pi_t = \beta\,\mathbb{E}_t\pi_{t+1} + \kappa^{mono}\,\tilde{y}_t
  \label{eq:NKPC_linearized}
\end{equation}
The HANK IS curve, which adds a redistribution channel $\Psi_t$ to the
standard Euler equation:
\begin{equation}
  \tilde{y}_t = \mathbb{E}_t\tilde{y}_{t+1} -
  \underbrace{\mathbf{G}^{C,r}[0,0]}_{\text{HANK: }{-0.137}}
  (i_t - \mathbb{E}_t\pi_{t+1} - \bar r)
  + \underbrace{\Psi_t}_{\substack{\text{redistribution} \\ \text{channel}}}
  \label{eq:IS_linearized}
\end{equation}
where $\Psi_t = (d\mathcal{M}_t/d\mu_t^w)\,\hat\mu_t^w$ captures the
novel channel: when monopsony deepens, the distribution of assets shifts
toward constrained high-MPC households, amplifying demand contractions.
The representative-agent benchmark has $\mathbf{G}^{C,r}[0,0] = -0.495$
(3.6$\times$ larger in magnitude), showing how HANK dampens the direct
intertemporal channel while leaving the wage income channel---and thus the
monopsony redistribution term---intact. Full market-clearing conditions and
the steady-state computation are in the Supplemental Appendix.
